\section{Introduction}
Let $M$ be a finitely generated graded module over a polynomial ring $R = \K[x_1,\ldots,x_n]$, where
$\K$ is a field. The Castelnuovo--Mumford regularity (or simply, regularity) $\reg(M)$ of $M$
is defined as
$$
 \reg(M)=\max \{j-i \mid \Tor_i^R(M, \K)_j \neq 0\}.
$$
 $\text{  }$ $\text{  }$Regularity  is an important invariant in commutative algebra and algebraic geometry that measures in some sense the complexity of ideals, modules, and sheaves. A question that has been studied by many is how the regularity behaves with respect to taking powers of homogeneous ideals. It is known that in the long-run $\reg(I^k)$ is linear in $k$, that is, there exist integers $a(I), b(I), c(I)$ such that $\reg(I^k) = a(I)k + b(I)$ for all $k\geq c(I)$ (see \cite{CHT,vijay}). For various classes of ideals people have studied these integers and also have looked for various upper and lower bounds for $\reg(I^k)$.
